\documentclass[10pt,twocolumn,letterpaper]{article}

\usepackage{cvpr}
\usepackage{float}
\usepackage{hyperref}

\graphicspath{{figures/}}
\hypersetup{
  colorlinks=true,
  linkcolor=blue,
  citecolor=blue,
  urlcolor=blue,
  pdftitle={Spatial Structure, Memorisation and Class Information in CIFAR-10},
  pdfauthor={Gourav}
}

% Generated by the existing build_report.py from the experiment results.
\newcommand{\Repo}{\url{https://github.com/Gourav-04/ML_assignment-3.git}}
\newcommand{\Commit}{\texttt{8b7703628a78dba75164b250316a7a396e84c7de}}
\newcommand{\LocalResult}{The clean patch network scored 52.10\% and random grouping scored 48.40\%. The signed clean-minus-shuffled difference was +3.7 points. The ordinary dense MLP scored 41.57\%. The paired patch difference cannot be attributed to a change in parameter count. Contiguous neighbourhoods helped the shared first layer.}
\newcommand{\MemoryParams}{1,707,274}
\newcommand{\MemoryResult}{The clean validation threshold was crossed at epoch 1. Assigned-label accuracy crossed at epoch 33. Under this definition, transferable learning preceded substantial random-label fitting. At epoch 65 the MLP predicted the assigned label on 61.20\% of flipped examples. Clean validation peaked at 38.74\% in epoch 30 and ended at 33.94\%. On the identical flipped images, original-label accuracy ended at 13.30\%.}
\newcommand{\SpatialResult}{The clean-minus-shuffled gap was +13.8 percentage points for the CNN and -1.6 for the dense MLP. Holding tiles intact changed CNN accuracy by +6.5 points relative to full shuffling. This pattern supports sensitivity to local neighbourhoods. Tile boundaries and changed global positions limit finer attribution.}
\newcommand{\ConfidenceResult}{CNN mean confidence on errors was 41.63\% on clean images, 31.89\% on pixel images, 37.43\% on tile images. Full shuffling changed the error count by +1384 out of 10,000 and mean error confidence by -9.7 points.}
\newcommand{\ReconstructionResult}{Validation MSE selected width 128; test reconstruction errors also appear for the representation comparison required in Q3(c).}
\newcommand{\DenoiseResult}{Frozen-CNN test accuracy was 26.78\% for corrupted, 26.28\% for clean AE, 27.05\% for denoising AE. Denoising changed accuracy by +0.27 percentage points relative to the corrupted image. It corrected 2379 previously wrong predictions and spoiled 2352 previously correct ones (paired exact McNemar $p=0.705$). The observed net change gives little evidence of a reliable improvement. The original clean-input CNN accuracy of 52.66\% shows how much discriminative information remains unrecovered.}
\newcommand{\ProbeResult}{In order of increasing test MSE the widths were 128, 32, 8. In descending probe accuracy the widths were 128, 32, 8. The rankings agree in this sweep. Width is a shared explanatory variable, so this does not establish that improving MSE causes better classification.}
\newcommand{\SvmResult}{Validation selected rbf with $C=1$; the frozen SVM achieved 42.79\% on test. The winner lies at the upper end of the evaluated grid; validation does not establish that it is a global optimum. For linear, validation peaked at $C=1$ (37.06\%); raising $C$ from 0.0001 to 1 moved training accuracy from 28.88\% to 45.14\%. At $C=1$ its train--validation gap was 8.08 points. For rbf, validation peaked at $C=1$ (42.10\%); raising $C$ from 0.0001 to 1 moved training accuracy from 33.06\% to 61.62\%. At $C=1$ its train--validation gap was 19.52 points. For poly, validation peaked at $C=1$ (33.64\%); raising $C$ from 0.0001 to 1 moved training accuracy from 11.46\% to 60.62\%. At $C=1$ its train--validation gap was 26.98 points. A wide positive gap at high $C$ indicates greater sensitivity to the training sample. The support-vector fraction in the plot counts each training point once across the binary subproblems.}
\newcommand{\PairResult}{The combined error counts were cat/dog (353), automobile/truck (352), deer/frog (321). Of the six near-boundary test examples, 4 received an incorrect multiclass prediction, including 3 predictions outside the inspected pair.}
\newcommand{\ClusterResult}{Validation silhouette by $k$ was 8 (0.077), 10 (0.067), 15 (0.058), 20 (0.044). The selected $k=8$ gave test ARI 0.049 and NMI 0.082. The cat and deer cluster histograms overlap by 90.40\% (sum of the smaller class-conditional mass in each cluster), while only 7.65\% of the two classes' images were confused with one another by the SVM. Similar unsupervised assignments can coexist with a separating supervised boundary.}
\newcommand{\Headline}{The CNN's clean-to-shuffled accuracy gap was 13.84 points; the dense MLP's gap was -1.64 points.}
\newcommand{\Repository}{\url{https://github.com/Gourav-04/ML_assignment-3.git}}
\newcommand{\Revision}{8b7703628a78dba75164b250316a7a396e84c7de}
\newcommand{\ChosenWidth}{128}
\newcommand{\CleanAccuracy}{52.66\%}
\newcommand{\PopulationRows}{Q1(a), Q2, Q3 & 45,000 & 4,500\\
Q1(b) memorisation & 4,000 & 400\\
Q4 SVM and clustering & 5,000 & 500\\}
\newcommand{\LocalRows}{Dense, clean & 101,098 & 42.26 & 41.57\\
Patch, contiguous & 100,410 & 51.94 & 52.10\\
Patch, shuffled & 100,410 & 49.48 & 48.40\\}
\newcommand{\MemoryRows}{1 & 22.00 & 18.20 & 9.60 & 20.88\\
15 & 46.03 & 34.40 & 12.70 & 35.66\\
30 & 59.27 & 32.30 & 18.20 & 38.74\\
33 & 57.73 & 31.10 & 20.40 & 37.54\\
50 & 71.87 & 23.00 & 38.30 & 34.92\\
65 & 84.63 & 13.30 & 61.20 & 33.94\\}
\newcommand{\SpatialRows}{MLP & clean & 101,098 & 41.57 & 5,843 & 34.50\\
MLP & pixel & 101,098 & 43.21 & 5,679 & 35.45\\
MLP & tile & 101,098 & 40.59 & 5,941 & 34.71\\
CNN & clean & 100,074 & 52.66 & 4,734 & 41.63\\
CNN & pixel & 100,074 & 38.82 & 6,118 & 31.89\\
CNN & tile & 100,074 & 45.28 & 5,472 & 37.43\\}
\newcommand{\EncoderRows}{8 & 80,283 & 0.02335 & 0.02346 & 29.22\\
32 & 117,171 & 0.01444 & 0.01450 & 34.53\\
128 & 264,723 & 0.01132 & 0.01136 & 39.78\\}
\newcommand{\DenoiseRows}{Clean image (Q2) & 52.66 & 4,734 & 41.63\\
corrupted & 26.78 & 7,322 & 54.83\\
clean AE & 26.28 & 7,372 & 37.92\\
denoising AE & 27.05 & 7,295 & 38.45\\}
\newcommand{\SweepRows}{linear & 0.0001 & 28.88 & 28.12 & 99.90\\
linear & 0.001 & 35.76 & 34.54 & 97.92\\
linear & 0.01 & 39.46 & 35.38 & 95.48\\
linear & 0.1 & 41.42 & 35.60 & 93.76\\
linear & 1 & 45.14 & 37.06 & 92.50\\
rbf & 0.0001 & 33.06 & 31.22 & 100.00\\
rbf & 0.001 & 33.06 & 31.22 & 100.00\\
rbf & 0.01 & 33.06 & 31.24 & 100.00\\
rbf & 0.1 & 38.24 & 35.42 & 99.10\\
rbf & 1 & 61.62 & 42.10 & 95.34\\
poly & 0.0001 & 11.46 & 11.08 & 100.00\\
poly & 0.001 & 11.46 & 11.08 & 100.00\\
poly & 0.01 & 15.84 & 14.80 & 99.78\\
poly & 0.1 & 33.36 & 27.50 & 98.26\\
poly & 1 & 60.62 & 33.64 & 95.36\\}
\newcommand{\PairRows}{cat, dog & 142 & 211 & 353 & 17.65\%\\
automobile, truck & 174 & 178 & 352 & 17.60\%\\
deer, frog & 189 & 132 & 321 & 16.05\%\\}
\newcommand{\ClusterRows}{0 & 70 & 214 & 69 & 97 & 80 & 70 & 63 & 209 & 160 & 400\\
1 & 415 & 109 & 131 & 53 & 68 & 88 & 24 & 74 & 447 & 144\\
2 & 16 & 90 & 89 & 120 & 95 & 70 & 155 & 48 & 15 & 26\\
3 & 252 & 72 & 91 & 78 & 24 & 62 & 41 & 35 & 109 & 62\\
4 & 69 & 88 & 162 & 136 & 136 & 140 & 145 & 181 & 48 & 125\\
5 & 110 & 235 & 177 & 180 & 205 & 137 & 168 & 113 & 153 & 130\\
6 & 22 & 93 & 138 & 120 & 167 & 131 & 258 & 123 & 8 & 65\\
7 & 46 & 99 & 143 & 216 & 225 & 302 & 146 & 217 & 60 & 48\\}
\newcommand{\LocalFindings}{The clean patch network scored 52.10\% and random grouping scored 48.40\%. The signed clean-minus-shuffled difference was +3.7 points. The ordinary dense MLP scored 41.57\%. The paired patch difference cannot be attributed to a change in parameter count. Contiguous neighbourhoods helped the shared first layer. The dense-to-patch difference also includes changes in connectivity and sharing.}
\newcommand{\MemoryFindings}{The clean validation threshold was crossed at epoch 1. Assigned-label accuracy crossed at epoch 33. Under this definition, transferable learning preceded substantial random-label fitting. At epoch 65 the MLP predicted the assigned label on 61.20\% of flipped examples. Clean validation peaked at 38.74\% in epoch 30 and ended at 33.94\%. On the identical flipped images, original-label accuracy ended at 13.30\%. The threshold describes substantial random-label fitting and cannot date the first memorised example.}
\newcommand{\CleanFindings}{The clean CNN attained 52.66\% on the official test images, and the dense reference attained 41.57\%. Validation selected epochs 6 and 7, respectively. This establishes the clean baseline for testing the locality hypothesis.}
\newcommand{\PixelFindings}{Pixel shuffling changed CNN accuracy by -13.84 points and MLP accuracy by +1.64 points. The fixed permutation preserves every input value while changing the neighbours presented to each convolution.}
\newcommand{\TileFindings}{The clean-minus-shuffled gap was +13.8 percentage points for the CNN and -1.6 for the dense MLP. Holding tiles intact changed CNN accuracy by +6.5 points relative to full shuffling. This pattern supports sensitivity to local neighbourhoods. Tile boundaries and changed global positions limit finer attribution.}
\newcommand{\ConfidenceFindings}{CNN mean confidence on errors was 41.63\% on clean images, 31.89\% on pixel images, 37.43\% on tile images. Full shuffling changed the error count by +1384 out of 10,000 and mean error confidence by -9.7 points. These conditional averages describe a different set of mistaken images in each condition.}
\newcommand{\EncoderSelection}{Validation MSE selected width 128; test reconstruction errors also appear for the representation comparison required in Q3(c).}
\newcommand{\ReconstructionFindings}{The selected width represents 3,072 image coordinates using 128 latent coordinates. This ratio compares dimensions; the code and the input image use different numerical precision.}
\newcommand{\DenoiseFindings}{Frozen-CNN test accuracy was 26.78\% for corrupted, 26.28\% for clean AE, 27.05\% for denoising AE. Denoising changed accuracy by +0.27 percentage points relative to the corrupted image. It corrected 2379 previously wrong predictions and spoiled 2352 previously correct ones (paired exact McNemar $p=0.705$). The observed net change gives little evidence of a reliable improvement. The original clean-input CNN accuracy of 52.66\% shows how much discriminative information remains unrecovered. In total, 4,731 images changed correctness. The net change conceals how often the classifier reversed individual decisions.}
\newcommand{\ProbeFindings}{In order of increasing test MSE the widths were 128, 32, 8. In descending probe accuracy the widths were 128, 32, 8. The rankings agree in this sweep. Width is a shared explanatory variable, so this does not establish that improving MSE causes better classification. The strongest frozen readout attained 39.78\%. The experiment measures information accessible to a linear classifier.}
\newcommand{\SweepFindings}{At small $C$, classification violations incur a relatively small penalty. Jointly low training and validation scores suggest underfitting. The table also shows when additional fitting on training data has a smaller return on validation.}
\newcommand{\SvmFindings}{Validation selected rbf with $C=1$; the frozen SVM achieved 42.79\% on test. The winner lies at the upper end of the evaluated grid; validation does not establish that it is a global optimum. For linear, validation peaked at $C=1$ (37.06\%); raising $C$ from 0.0001 to 1 moved training accuracy from 28.88\% to 45.14\%. At $C=1$ its train--validation gap was 8.08 points. For rbf, validation peaked at $C=1$ (42.10\%); raising $C$ from 0.0001 to 1 moved training accuracy from 33.06\% to 61.62\%. At $C=1$ its train--validation gap was 19.52 points. For poly, validation peaked at $C=1$ (33.64\%); raising $C$ from 0.0001 to 1 moved training accuracy from 11.46\% to 60.62\%. At $C=1$ its train--validation gap was 26.98 points. A wide positive gap at high $C$ indicates greater sensitivity to the training sample. The support-vector fraction in the plot counts each training point once across the binary subproblems.}
\newcommand{\PairExamples}{The combined error counts were cat/dog (353), automobile/truck (352), deer/frog (321). Of the six near-boundary test examples, 4 received an incorrect multiclass prediction, including 3 predictions outside the inspected pair. Test image 2233 is a cat with score -0.000942 and predicted class airplane. Test image 9652 is a dog with score -0.000475 and predicted class ship. Test image 8366 is a automobile with score +0.000617 and predicted class automobile. Test image 2482 is a truck with score +0.000056 and predicted class automobile. Test image 6154 is a deer with score +0.005962 and predicted class airplane. Test image 1245 is a frog with score -0.000017 and predicted class frog.}
\newcommand{\ClusterSelection}{K-means uses ten initializations per candidate on the training subset. Validation silhouette on the fixed sample of 2,000 by candidate count was 8 (0.0771), 10 (0.0673), 15 (0.0576), 20 (0.0439). The selected count was 8. We assigned official test images to centroids before inspecting their labels.}
\newcommand{\ClusterFindings}{Test ARI was 0.0488 and test NMI was 0.0817. ARI compares pairwise assignment agreements after adjustment for chance. NMI measures the information shared by cluster identity and class.}
\newcommand{\OverlapFindings}{The class distributions for cat and deer overlap by 90.40\% across clusters. The SVM directly swaps these two classes on 7.65\% of their test examples. Overlap sums the smaller class-conditional mass in each cluster; the swap rate counts supervised errors into the other class. Similar coarse cluster memberships can coexist with a separating supervised boundary.}
\newcommand{\AbstractFindings}{Which image properties support recognition on CIFAR10? We intervene on neighbourhood structure, labels and reconstruction targets in one stratified split. The CNN's clean-to-shuffled accuracy gap was 13.84 points; the dense MLP's gap was -1.64 points. Denoising corrected 2,379 previously incorrect predictions and introduced 2,352 errors. On a frozen autoencoder representation the validation-selected SVM attained 42.79\% test accuracy. The paired comparisons distinguish pixel fidelity, linear class readout and unsupervised cluster membership.}
\newcommand{\Conclusions}{Preserving intact tiles changed CNN accuracy by +6.46 points relative to full pixel shuffling. In order of increasing test MSE the widths were 128, 32, 8. In descending probe accuracy the widths were 128, 32, 8. The rankings agree in this sweep. Width is a shared explanatory variable, so this does not establish that improving MSE causes better classification. The paired counts in the denoising experiment yield a net change of +27 successful predictions. These results concern the selected readouts and fixed training budgets. The supervised and unsupervised class-pair comparison further separates global cluster membership from decision-boundary geometry.}


% Optional on the first build. Step 4 creates this file after the source push.
\InputIfFileExists{source_commit.tex}{}{}

\newcommand{\ReportSourceLink}{%
  \ifdefined\ReportCommit
    \href{https://github.com/Gourav-04/ML_assignment-3/commit/\ReportCommit}%
         {\texttt{\ReportShortCommit}}%
  \else
    Pending source commit%
  \fi
}

\title{Spatial Structure, Memorisation and Class Information in CIFAR-10}

\author{Gourav\\
Machine Learning, Assignment 3\\
Roll number: \texttt{m26cse018}\\
{\small Repository:
\href{https://github.com/Gourav-04/ML_assignment-3}
     {\texttt{Gourav-04/ML\_assignment-3}}}\\
{\small Source commit: \ReportSourceLink}}

\newcommand{\fig}[4]{%
  \begin{figure}[H]
    \centering
    \includegraphics[width=#1\linewidth]{#2}
    \caption{#3}
    \label{#4}
  \end{figure}
}

\newcommand{\tbl}[3]{%
  \begin{table}[H]
    \centering
    \caption{#1}
    \label{#2}
    \small
    #3
  \end{table}
}

\begin{document}
\maketitle
\raggedbottom

\begin{abstract}
\AbstractFindings{}
\end{abstract}

\section{Dataset and experimental design}

CIFAR-10 contains $32\times32$ colour images from ten
classes~\cite{krizhevsky2009}. We create one stratified split from the
50,000 official training images. Reading the digits of
\texttt{m26cse018} in order gives seed 26018. Each class contributes
4,500 training images and 500 validation images. The official test set
contains a further 1,000 images per class. We scale RGB values to
$[0,1]$ and initialise the networks from random weights.

\paragraph{Model selection.}
Classifiers minimise cross-entropy. We keep the earliest checkpoint
with the highest validation accuracy. Autoencoders minimise mean
squared error (MSE), and we keep the earliest checkpoint with the
lowest validation MSE. Adam uses a learning rate of $10^{-3}$ and
batches of 256. We measure training and validation curves in evaluation
mode after each epoch. Q1(a), Q2 and the autoencoders use seven epochs.
Q1(b) uses a larger MLP and 65 epochs to track the fitting of
incorrect labels. We select the SVM and clustering configurations
using validation data before measuring official test performance.

\paragraph{Comparisons across questions.}
Q1(a) changes which pixels enter a patch while holding the patch
architecture fixed. Q2 reuses the dense MLP and compares its response
to spatial transformations with that of a CNN. Q3(b) freezes the clean
CNN from Q2(a), so each reconstruction method is evaluated with the
same classifier. Q3(c) measures how well a linear probe can read class
information from a frozen encoder. Q4 uses the clean encoder selected
in Q3(a) for both classification and clustering.

\tbl{Training populations. Both smaller subsets are drawn from the
training split before model selection.}{tab:units}{
\begin{tabular}{lrr}
\toprule
Experiment & Images & Per class\\
\midrule
\PopulationRows
\bottomrule
\end{tabular}}

\section{Q1. Local grouping and memorisation}

\subsection{Q1(a). Does grouping neighbours help an MLP?}

A dense first layer can use any input pixel, but its connectivity
contains no explicit information about which pixels are neighbours.
To test whether meaningful local groups help, we divide each image
into 64 patches. We flatten each $4\times4$ RGB patch into 48 values
and apply the same $48\to32$ linear layer and ReLU to every patch.
The resulting $64\times32=2{,}048$ activations feed a 48-unit ReLU
layer and a ten-class output layer. Concatenation preserves the
position of each patch.

We train the patch model twice. The first run receives ordinary
images. In the second, one fixed permutation rearranges all 1,024
RGB triplets before patches are formed. Both runs see the same pixel
values and use the same number of parameters, connections, optimiser
steps and epochs. They also share the initialisation seed and data
order. Their difference therefore measures the effect of grouping
neighbours under this architecture. A conventional dense MLP with
hidden widths 32 and 64 provides a reference.

\tbl{Q1(a) model comparison. Validation selects each checkpoint.
Accuracy values are percentages.}{tab:q1}{
\begin{tabular}{lrrr}
\toprule
Model / grouping & Parameters & Val. & Test\\
\midrule
\LocalRows
\bottomrule
\end{tabular}}

\fig{1}{q1_curves.pdf}
{Q1(a) post-epoch accuracy. Dashed lines show training accuracy and
solid lines show validation accuracy. Both patch conditions use the
same update budget.}{fig:q1}

\LocalFindings{}

\subsection{Q1(b). When does the MLP fit deliberately wrong labels?}

We train a separate $3072\to512\to256\to10$ ReLU MLP with 1,707,274
parameters. Within a fixed subset of 4,000 training images, we assign
1,000 images an incorrect label. Adding a random integer from 1 to 9
modulo 10 ensures that the assigned label differs from the original.
Each image keeps its assigned label throughout the 65 epochs.
Validation images retain their original labels.

For the flipped images, we evaluate each prediction against both the
original and the assigned label. For example, if a cat is assigned
the label truck, predicting cat counts toward original-label accuracy
and predicting truck counts toward assigned-label accuracy. Predicting
deer counts toward neither. Clean validation accuracy measures whether
the network has learned patterns that transfer to unseen images.

\fig{1}{memory.pdf}
{Q1(b) accuracy on unchanged training labels, original and assigned
labels for flipped images, and clean validation images. The two
flipped-label curves use exactly the same 1,000 images.}{fig:mem}

\tbl{Q1(b) trajectory at selected epochs. Original and assigned refer
to the flipped images. All values are percentages.}{tab:mem}{
\begin{tabular}{rrrrr}
\toprule
Epoch & Unchanged & Original & Assigned & Val.\\
\midrule
\MemoryRows
\bottomrule
\end{tabular}}

\MemoryFindings{}

\section{Q2. Which spatial relationships matter to a CNN?}

\subsection{Q2(a). A budget-matched clean baseline}

The CNN uses three padded $3\times3$ convolutions with 32, 64 and 96
output channels and a ReLU after each convolution. The first
convolution uses stride two. The first two blocks each include
$2\times2$ max pooling. Adaptive average pooling produces a $2\times2$
grid of 96 channels. A $384\to64\to10$ head produces the logits.

Our comparison permits at most a 5\% parameter difference and fixes
the examples, epochs, optimiser steps and checkpoint-selection rule.
The CNN has 100,074 parameters and the dense MLP has 101,098, a
difference of about 1\%. Computational cost per update can still differ.

\CleanFindings{}

\subsection{Q2(b). Fixed permutation of individual pixels}

We apply one random bijection of spatial positions to every image,
keeping the three channels of each pixel together. A dense layer can
absorb a fixed input permutation into the columns of its weight
matrix. A convolutional kernel, however, now combines pixels that
could have been far apart in the original image. We therefore predict
a larger clean-to-permuted accuracy drop for the CNN. We retrain each
architecture on the permuted training images and evaluate it using
the same permutation.

\fig{.88}{interventions.pdf}
{Official test images 0 and 1 under the three Q2 conditions. Each row
contains the same RGB triplets. The last column preserves neighbourhoods
within individual tiles.}{fig:intervention}

\tbl{Q2 official test performance. Conf.\ is the mean maximum softmax
confidence among incorrect predictions. Accuracy and confidence are
percentages.}{tab:spatial}{
\setlength{\tabcolsep}{3pt}
\begin{tabular}{llrrrr}
\toprule
Model & Input & Params & Acc. & Errors & Conf.\\
\midrule
\SpatialRows
\bottomrule
\end{tabular}}

\PixelFindings{}

\subsection{Q2(c). Preserving neighbourhoods inside tiles}

To distinguish local neighbourhoods from the arrangement of larger
parts, we shuffle 64 intact $4\times4$ tiles using one fixed map.
Pixels within a tile keep their original neighbours, while the layout
between tiles changes (Fig.~\ref{fig:intervention}). We train fresh
MLP and CNN checkpoints on these transformed images.

\fig{1}{q2_curves.pdf}
{Q2 training and validation accuracy. Each curve uses a separately
trained model under its corresponding input condition.}{fig:q2}

\TileFindings{}

\ConfidenceFindings{}

\section{Q3. Pixel fidelity and class information}

\subsection{Q3(a). Select a clean encoder by reconstruction error}

The encoder maps an image $x$ to a code $z=E(x)$, and the decoder
produces $\hat{x}=D(z)$. Training uses the image itself as the target.
Three $4\times4$ convolutions with stride two and padding one produce
16, 32 and 48 channels. The spatial resolution falls from $32\times32$
to $4\times4$. A linear layer maps the resulting 768 activations to
8, 32 or 128 coordinates.

The decoder expands the code to 768 activations and uses three mirrored
transposed convolutions. Hidden layers use ReLU and the final sigmoid
restricts reconstructed values to $[0,1]$. Varying the bottleneck
changes the two weight matrices touching the code. The convolutional
blocks and training settings are held fixed.

For $N$ images, the reconstruction error is
\begin{equation}
\mathrm{MSE} =
\frac{1}{3072N}\sum_{i=1}^{N}\|x_i-D(E(x_i))\|_2^2 .
\end{equation}

\EncoderSelection{}

\fig{1}{ae_learning.pdf}
{Validation reconstruction MSE from the epoch logs of the three clean
autoencoders.}{fig:aecurve}

\tbl{Q3(a,c) reconstruction error and class readout. The probe column
gives official test accuracy.}{tab:ae}{
\begin{tabular}{rrrrr}
\toprule
Width & Parameters & Val.\ MSE & Test MSE & Probe (\%)\\
\midrule
\EncoderRows
\bottomrule
\end{tabular}}

\ReconstructionFindings{}

\subsection{Q3(b). Can a denoiser restore the CNN's decisions?}

The denoising objective reconstructs a clean target from a corrupted
input~\cite{vincent2010}. We add independent Gaussian noise with
standard deviation 0.2 to the scaled RGB components and clip to
$[0,1]$. A new width-\ChosenWidth{} model minimises MSE from each
corrupted image to its original clean target. We regenerate training
noise across epochs and keep the evaluation noise draw fixed for
each image. The ordinary width-\ChosenWidth{} autoencoder from
Q3(a) provides a second reconstruction baseline.

We freeze the clean CNN from Q2(a) and give it three views of each
test example: the noisy image, the ordinary autoencoder's reconstruction
of that noisy image, and the denoising autoencoder's reconstruction.
The clean-image accuracy of \CleanAccuracy{} provides a reference for
the loss under corruption.

\tbl{Q3(b) results from the same frozen CNN. Confidence is averaged
over incorrect predictions.}{tab:denoise}{
\begin{tabular}{lrrr}
\toprule
CNN input & Acc.\ (\%) & Errors & Conf.\ (\%)\\
\midrule
\DenoiseRows
\bottomrule
\end{tabular}}

\fig{.96}{denoising_examples.pdf}
{First four official test images in file order. Each row gives the
index and true class. The label under each view is the frozen CNN's
prediction. Both reconstructions receive the same noisy input.}
{fig:denoise}

Reading across a row of Fig.~\ref{fig:denoise} shows which individual
decisions change. The paired counts use the entire official test set
and quantify both corrected predictions and newly introduced errors.

\DenoiseFindings{}

\subsection{Q3(c). What can a linear classifier read from each code?}

We freeze each clean encoder, extract its latent coordinates, and
standardise each feature using training statistics. A multinomial
logistic regression uses $C=1$ and at most 1,500 iterations. At test
time we reuse the scaler and probe fitted for that width. The encoder
receives no gradient from the supervised readout. This measures the
class information accessible to this particular linear classifier.

\ProbeFindings{}

\section{Q4. Margins and clusters in one representation}

\subsection{Q4(a). Selecting the SVM kernel and penalty}

We fit support-vector classifiers~\cite{cortes1995} to representations
from the frozen width-\ChosenWidth{} encoder selected in Q3(a).
A stratified 5,000-image training subset supplies the feature scaler,
every SVM fit and every clustering fit. All 5,000 validation images
are available for SVM selection. We apply the training scaler to
validation and test data.

We compare linear, Gaussian RBF and cubic polynomial kernels.
The nonlinear kernels use scale-based gamma and the polynomial offset
is zero. Five $C$ values from $10^{-4}$ through $1$ span four orders
of magnitude. A preliminary range extending to 100 was abandoned before
official test evaluation because the large-$C$ linear solve was slow.
All kernels use the same final grid. We select the highest validation
accuracy, breaking ties by the configured sweep order. The support-vector
fraction is the number of distinct training support vectors divided
by 5,000.

\tbl{Q4(a) kernel sweep. Training accuracy, validation accuracy and
support-vector fraction are percentages.}{tab:svm}{
\begin{tabular}{lrrrr}
\toprule
Kernel & $C$ & Train & Val. & SV\\
\midrule
\SweepRows
\bottomrule
\end{tabular}}

\fig{1}{svm.pdf}
{Validation accuracy and support-vector fraction across the penalty
grid. The horizontal axis is logarithmic. The curves show fractions
and Table~\ref{tab:svm} gives percentages.}{fig:svm}

\SweepFindings{}

\SvmFindings{}

\subsection{Q4(b). Difficult pairs and the decision boundary}

We compute the test confusion matrix after selecting the SVM. Rows
are true classes and columns are predictions. We rank each unordered
pair $(i,j)$ by $M_{ij}+M_{ji}$. Since each class has 1,000 test images,
this ranking also orders the pairwise swap rates. We use the resulting
pairs to describe the selected model's errors.

\fig{1}{svm_confusion.pdf}
{Selected SVM test confusion matrix, normalised by true-class count.
Off-diagonal cells identify errors. Both directions contribute to the
ranking of difficult pairs.}{fig:confusion}

\tbl{Direct errors for the three most confused pairs. Each rate uses
the 2,000 test images belonging to that pair.}{tab:hard}{
\begin{tabular}{lrrrr}
\toprule
Pair $(i,j)$ & $i\to j$ & $j\to i$ & Total & Rate\\
\midrule
\PairRows
\bottomrule
\end{tabular}}

For each pair we examine its one-versus-one decision function. Among
training support vectors with a nonzero dual coefficient for that
binary decision, we choose the smallest absolute score in each true
class. We make an analogous selection among test images from the two
classes. An absolute score near zero places an example near the binary
decision boundary in score terms. The sign identifies the favoured
class; positive signs favour the first class in CIFAR-10 label order.

\fig{1}{pairs.pdf}
{Rows follow Table~\ref{tab:hard}. The first two columns show training
support vectors. The last two show near-boundary test images. Each
label includes its signed pair score.}{fig:pairs}

\PairExamples{}

A near tie for one binary classifier does not determine the ten-way
vote. Images can also have similar decision scores while differing
visually. Figure~\ref{fig:pairs} illustrates selected boundary cases;
the confusion matrix gives the population counts.

\subsection{Q4(c). What do unlabeled clusters recover?}

\ClusterSelection{}

\tbl{Q4(c) test contingency table. Each column sums to 1,000. Each row
corresponds to one selected cluster.}{tab:cluster}{
\resizebox{\linewidth}{!}{%
\begin{tabular}{r|rrrrrrrrrr}
\toprule
Cluster & Air & Car & Bird & Cat & Deer & Dog & Frog & Horse & Ship & Truck\\
\midrule
\ClusterRows
\bottomrule
\end{tabular}}\\[-1pt]
{\footnotesize Air = airplane. Car = automobile.}}

\ClusterFindings{}

\OverlapFindings{}

\section{Limits and reproducibility}

The comparisons apply to this split, these architectures and these
training budgets. Another seed or longer training could change the
clean accuracies and intervention effects. Matching parameters and
optimiser steps still permits different computational costs. Confidence
on errors describes the subset of examples each model misclassifies.
The Q3 probes measure information accessible to a linear readout, and
the Q4 clustering comparison covers the tested $K$-means candidates.

The project saves split and subset indices, flipped-label assignments,
epoch logs, selected weights, the kernel sweep and the final tables.
The four question scripts implement the individual experiments.
\texttt{run\_all.py} runs the assignment in question order. With saved
checkpoints it reuses neural training trajectories and recomputes the
downstream fits and test measurements. The intervention and denoising
figures show the first test examples in their original order. Reported
test accuracies use the complete official test set.

\paragraph{Reproducing the report.}
From the project directory, run \texttt{python run\_all.py} to reproduce
the experiments. With the saved results, checkpoints and dataset
available, run \texttt{python build\_report.py} to regenerate the
tables, figures and PDF. The files in Table~\ref{tab:audit} allow
the rounded values in the report to be checked against the recorded
measurements.

\tbl{Evidence files in the project's results directory.}{tab:audit}{
\begin{tabular}{ll}
\toprule
Part & Evidence file\\
\midrule
Q1(a) & \texttt{q1\_test.csv}\\
Q1(b) & \texttt{memory.csv}\\
Q2(a--c) & \texttt{q2\_test.csv}\\
Q3(a,c) & \texttt{representation.csv}\\
Q3(b) & \texttt{denoising.csv}\\
Q3(b), paired counts & \texttt{paired\_denoising.json}\\
Q4(a) & \texttt{svm\_sweep.csv}\\
Q4(b) & \texttt{pairs.csv}\\
Q4(c) & \texttt{contingency.csv}\\
Q4 final metrics & \texttt{bonus.json}\\
\bottomrule
\end{tabular}}

\section{Conclusion}
\Conclusions{}

\begin{thebibliography}{9}

\bibitem{krizhevsky2009}
Alex Krizhevsky.
\newblock Learning Multiple Layers of Features from Tiny Images.
\newblock Technical report, University of Toronto, 2009.
\newblock
\href{https://www.cs.toronto.edu/~kriz/learning-features-2009-TR.pdf}
{Technical report}.

\bibitem{vincent2010}
Pascal Vincent, Hugo Larochelle, Isabelle Lajoie, Yoshua Bengio,
and Pierre-Antoine Manzagol.
\newblock Stacked Denoising Autoencoders: Learning Useful Representations
in a Deep Network with a Local Denoising Criterion.
\newblock \emph{Journal of Machine Learning Research}, 11:3371--3408, 2010.
\newblock
\href{https://www.jmlr.org/papers/v11/vincent10a.html}{Article}.

\bibitem{cortes1995}
Corinna Cortes and Vladimir Vapnik.
\newblock Support-vector networks.
\newblock \emph{Machine Learning}, 20:273--297, 1995.
\newblock
\href{https://doi.org/10.1007/BF00994018}{DOI}.

\end{thebibliography}

\end{document}